\section{Very Special Problems}

In this question, we consider special cases of NP-complete problems. Formally, we say that Problem B is a \textbf{special case} of Problem A if every instance of Problem B can be viewed as an instance of Problem A. We've seen examples of this in class: for example, 3-SAT is a special case of SAT (where every clause has length 3). Minimum spanning tree is a special case of Steiner Tree, in which the set of vertices we need to connect includes all the vertices in $V$.

\begin{questions}

\question[4] Consider the \textbf{Bounded-Leaf Spanning Tree Problem (BLST)}: given a graph $G = (V,E)$ and an integer $k$, does $G$ have a spanning tree with no more than $k$ leaves?

Give an NP-complete problem that is a special case of BLST, and justify why this problem is a special case.

\ifsolutions\input{q3a-sol.tex}\fi

\question[3] You showed in the previous question that there is an NP-complete problem that is a special case of BLST, and it is not difficult to show that BLST is in NP (though we are not asking you to do this). Does this imply that BLST is NP-complete? Justify your answer.

\ifsolutions\input{q3b-sol.tex}\fi

\question[4] Give an example of a polytime-solvable problem you have seen in this class that is a special case of an NP-complete problem. Justify your answer.

\ifsolutions\input{q3c-sol.tex}\fi

\end{questions}